\FloatBarrier
\section{Declaración de uso de IA}\label{sec:ia}

En este proyecto usamos asistentes de programación con IA para escribir código, pruebas y documentos. Durante la revisión del 25 de noviembre se utilizó ChatGPT (GPT-6.2) como apoyo para interpretar evidencias de red, redactar propuestas de documentación y actualizar archivos en una rama independiente mediante GitHub. La integrante revisó las propuestas y obtuvo las evidencias reales en el navegador y Excel.

Usé ChatGPT (GPT-6.2) como apoyo para integrar mis aportaciones al proyecto, contrastar la información con la documentación existente y revisar la redacción. También lo utilicé para identificar inconsistencias y proponer correcciones en mi parte de la entrega.

\pend{declaraciones de herramientas, modelos y usos de los demás integrantes}

En la revisión final, después de la entrega 3, un agente escribió los planes de \tt{plans/} y otros agentes los ejecutaron, como explica la \sref{sec:agentes}. El borrador de este reporte también lo escribió un agente, con un modelo Claude Opus de Anthropic, a partir de los documentos del repositorio. El índice de planes registra Sonnet para los planes 009 a 011 y Opus para el plan 012.

Los cambios se registraron mediante pull requests. Los PR 51 y 53 se fusionaron con ayuda de Codex, a petición de una integrante, después de comprobar que habían pasado CI. Los datos de este reporte salen del repositorio y de GitHub, y lo que no se pudo comprobar está marcado como pendiente. Las conclusiones propuestas y las lecciones individuales incorporadas se redactaron con apoyo de IA. La propuesta de conclusiones requiere aprobación grupal. El resto de las lecciones está pendiente.
